\documentclass{article}
\usepackage[utf8]{inputenc}
\usepackage{xcolor}
\usepackage{hyperref}

\usepackage{rotating}

\usepackage{tabls}
\usepackage{booktabs}
\usepackage{amsmath}
\usepackage{amssymb}
\usepackage{amsthm}
\usepackage{amsfonts}
\usepackage{multicol}
\usepackage{enumitem}
\usepackage{microtype}
\usepackage{tikz}
\usepackage{pgfplots}
\usetikzlibrary{positioning}
\usepackage{multirow}

\usepackage{comment}

\usepackage{graphicx}
\usepackage{wrapfig}

\theoremstyle{plain}
\newtheorem{theorem}{Theorem}
\newtheorem*{theorem*}{Theorem}
\newtheorem{lemma}[theorem]{Lemma}
\newtheorem*{lemma*}{Lemma}
\newtheorem{proposition}[theorem]{Proposition}
\newtheorem{corollary}[theorem]{Corollary}

\theoremstyle{box}
\newtheorem{definition}[theorem]{Definition}
\newtheorem*{axiom*}{Axiom}
\newtheorem{dfn}[theorem]{Definition}
\newtheorem*{definition*}{Definition}
\newtheorem*{exercise*}{Exercise}
\newtheorem{observation}[theorem]{Observation}
\newtheorem{remark}[theorem]{Remark}
\newtheorem{example}[theorem]{Example}
\newtheorem{question}[theorem]{Question}
\newtheorem*{prep-problems}{Preparation Problems}

\newcommand{\pd}[3]{\frac{\partial^{#3} {#1}}{\partial {#2}^{#3}}}
\newcommand{\fpd}[2]{\frac{\partial {#1}}{\partial {#2}}}
\newcommand{\diff}[2]{\frac{d {#1}}{d {#2}}}


\begin{document}
\section{Objectives}
\begin{itemize}
\item Calculate basic discrete probabilities.
	\begin{itemize}
	\item Compute the probability of equally likely events.
	\item Compute the probability of the intersection of two events.
	\item Compute the probability of the union of two events.
	\item Compute conditional probabilities.
	\item Recognize independent events.
	\item Recognize mutually exclusive events.
	\end{itemize}
\item Write down the pmf for some discrete random variables.
\item Use the pmf to calculate probability.
\end{itemize}
\subsection{Group Collaborative Spaces}
\begin{itemize}
\item Section 01
\begin{itemize}
\item \href{https://docs.google.com/document/d/1qphjKMfvf8ZYyWgxmPa8aR2uxxuY0OsACZgBQU16K-M/edit?usp=sharing}{Group 1 Collaboration Space}
\item \href{https://docs.google.com/document/d/1fO_BrHFA-t3y8tcKkPVKwctZH11N1Q1KoPfUf_HA55E/edit?usp=sharing}{Group 2 Collaboration Space}
\item \href{https://docs.google.com/document/d/1kP8tFnS9bRtFdNE6igO5XWp8jiEyH2SZUNBEGfZ5iis/edit?usp=sharing}{Group 3 Collaboration Space}
\end{itemize}
\item Section 02
\begin{itemize}
\item \href{https://docs.google.com/document/d/12h8t5T-amdit0sv4IV7YdSRlPrWIpbOd1DPqB-ygvQA/edit?usp=sharing}{Group 1 Collaboration Space}
\item \href{https://docs.google.com/document/d/1wzOKAlH0lwUrYuI6OBgbbvopKUznyT5rJF9nrlCT7gg/edit?usp=sharing}{Group 2 Collaboration Space}
\end{itemize}
\end{itemize}



%\subsection{table of contents}

\newpage
\section{Agenda}
\subsection{Probability}

\subsubsection{Examples}
Consider an experiment where we roll a 6-sided die once. The \textbf{sample space} in this case is $\Omega = \{1,2,3,4,5,6\}$.
\begin{itemize}
\item Let $A$ be the event the outcome is odd and $B$ be the event the outcome is greater than 3.\\
Thus $A = \{1,3,5\}$ and $B = \{4,5,6\}$. Calculate the following:
\begin{itemize}
\item $P(A)$
\item $P(B)$
\item $P(A \cap B)$
\item $P(A \cup B)$
\item $P(A|B)$
\item $P(B|A)$
\end{itemize}

\item Let $A$ be the event the outcome is less than 5 and $B$ be the event the outcome is even.\\
Thus $A = \{1,2,3,4\}$ and $B = \{2,4,6\}$. Calculate the following:
\begin{itemize}
\item $P(A)$
\item $P(B)$
\item $P(A \cap B)$
\item $P(A \cup B)$
\item $P(A|B)$
\item $P(B|A)$
\begin{definition*}[Disjoint Sets]
If the intersection of two sets is the empty set, we say the sets are disjoint.
\end{definition*}
\end{itemize}

\item Let $A$ be the event the outcome is odd and $B$ be the event the outcome is even.\\
Thus $A = \{1,3,5\}$ and $B = \{2,4,,6\}$. Calculate the following:
\begin{itemize}
\item $P(A)$
\item $P(B)$
\item $P(A \cap B)$
\item $P(A \cup B)$
\item $P(A|B)$
\item $P(B|A)$
\end{itemize}
\end{itemize}



\newpage
\subsection{Practice with Definitions}
The following definitions are from Section 3.1 and 3.2 of the OpenStax Introductory Statistics textbook.
\begin{definition*}[Probability]
The probability of any outcome is the long-term relative frequency of that outcome.
\end{definition*}

\begin{definition*}[Independent Events]
Two events A and B are independent if the knowledge that one occurred does not affect the chance the other occurs. Two events are independent if the following are true: $P(A|B) = P(A)$ and $P(B|A) = P(B)$. When two events are independent $P(A \cap B) = P(A)P(B)$.
\end{definition*}

\begin{definition*}[Mutually Exclusive Events]
A and B are mutually exclusive events if they cannot occur at the same time. This means that A and B do not share any outcomes and $P(A \cap B) = 0.$
\end{definition*}

\begin{exercise*}[Example from Prep]
The distribution of a random sample $S$ of 100 individuals, organized by gender and whether they are right-handed or left-handed is given in the table below.
\begin{itemize}
\item $M$ is the event that the subject is male
\item $F$ is the event that the subject is female
\item $R$ is the event that the subject is right-handed
\item $L$ is the event that the subject is left-handed
\end{itemize}

\begin{table}[h]
\begin{center}
\begin{tabular}{|c|c|c|}
\hline
& Right-handed & Left-handed\\
\hline
Males & 43 & 9\\
\hline
Females & 44 & 4\\
\hline
\end{tabular}
\end{center}
\end{table}
\end{exercise*}



\newpage
\subsection{Random Variables and Distributions}
\noindent The following definition is from Chapter 4 Introduction in the OpenStax Introductory Statistics textbook.
\begin{definition*}[Random Variable]
A random variable describes the outcomes of a statistical experiment in words.
\end{definition*}
Random variables have both a set of values they can take on and a probability that they take on each value. This information is collected in a distribution function (sometimes called a probability model) for the random variable. When the random variable is discrete, meaning the random variable has only a finite (or countably infinite) number of possible values, then the distribution function function is also called a probability mass function (pmf).

\begin{exercise*}[Example]
Let $R$ be the random variable that represents that value of the roll of a 6-sided die. Thus $R$ could be 1 or $R$ could be 2 or $R$ could be 3 or $R$ could be 4 or $R$ could be 5 or $R$ could be 6. The pmf for $R$ is 
      $$p(r) = 
    \begin{cases}
    \frac{1}{6} & r=1,2,3,4,5,6 \\
    0 & \text{otherwise}
    \end{cases}.$$

\begin{itemize}
\item What is the domain of $p$?
\item What is the range of $p$?
\item Compute $P(R = 3)$.
\item Calculate $P(R \neq 6)$.
\item Find $P(R > 4)$.
\item Given $A$ is the event that $R$ is an odd, calculate $P(A)$.
\end{itemize}
\end{exercise*}



\newpage
\subsection{PREP for Friday}
\begin{exercise*}[PMF and Probability]
Consider the experiment where two 6-sided dice are rolled and we sum the dice. Let $S$ be the random variable that represents the value of the sum of the dice.
\item Find the pmf for $S$.
\item Compute the following:
\begin{itemize}
\item $P(S = 2)$
\item $P(S = 8)$
\item $P(S = 10)$
\item $P(S \leq 3)$
\item $P(S \leq 9)$
\end{itemize}
\end{exercise*}
 
\end{document}